\chapter{Experiments}
\label{chap:experiments}

\section{Datasets}

Describe each evaluated target pack, the number of paired textures, the validation texture names, and any preprocessing differences such as original resolution.
Keep this section reproducible by recording the exact command used to create \texttt{data/processed/}.

\begin{table}[htbp]
  \centering
  \caption{Dataset summary template.}
  \begin{tabular}{lrrl}
    \toprule
    Target pack & Train pairs & Validation pairs & Notes \\
    \midrule
    \todo{pack-id} & \todo{n} & \todo{n} & \todo{resolution or filtering notes} \\
    \bottomrule
  \end{tabular}
\end{table}

\section{Training Configuration}

Report the model versions, optimizer settings, number of steps, learning-rate schedule, random seed policy, and hardware.
For generated results, also report the checkpoint path and inference command.

\begin{lstlisting}[language=bash,caption={Example commands to replace with the final experimental run.}]
uv run spritecraft preprocess
uv run spritecraft train --run report-baseline
uv run spritecraft train-recolor --run report-baseline
uv run spritecraft generate --run report-baseline
\end{lstlisting}

\section{Quantitative Metrics}

Use metrics to identify broad trends, not as the only evidence of quality.
For this project, candidate scores, reconstruction losses, color statistics, and sharpness penalties are useful diagnostics, but the final report should pair them with validation matrices and side-by-side texture comparisons.

\section{Qualitative Evaluation}

\placeholderfigure{Validation matrix placeholder}{Replace this placeholder with a grid from the checkpoint validation directory or an output comparison image.}

When reviewing images, comment on material identity, palette match, edge sharpness, contrast, and tiling behavior.
Record failure cases explicitly, especially for transparent textures, line-art textures, leaves, glass, and high-frequency blocks.

\section{Ablations}

Suggested ablations include removing the recolor fallback, changing the number of diffusion candidates, disabling detail injection, and varying the weights of content-aware losses.
Each ablation should include representative images as well as metric summaries.
